\chapter{INTRODUCTION} \label{c1}
\justifying
\section{Introduction}
Attendance management is a routine but critical academic process in all educational institutions. The quality of attendance records directly affects academic monitoring, internal assessments, and student discipline. In many classrooms, attendance is still recorded manually, which consumes valuable teaching time and introduces avoidable errors. Manual systems also make proxy attendance easier, where one student marks attendance for another.

In response to these problems, smart attendance systems have become an active area of implementation in campus environments. The objective is to automate attendance marking while ensuring both identity verification and physical presence verification. A secure system should be easy to use for teachers and students while preventing invalid submissions.

\section{Need for the Project}
The need for this project arises from three major concerns:
\begin{enumerate}
\item \textbf{Time efficiency in classrooms:} Manual attendance consumes the first few minutes of every class.
\item \textbf{Proxy attendance prevention:} OTP alone or manual roll-call cannot fully prevent impersonation.
\item \textbf{Reliable record keeping:} Institutions require structured, queryable, and auditable attendance records.
\end{enumerate}

The Smart Attendance System addresses these concerns by combining session OTP validation, Bluetooth proximity checks, and face verification before final attendance submission.

\section{Project Title}
\textbf{Smart Attendance System using Bluetooth and Verification}

\section{Project Objectives}
The primary objectives of this project are:
\begin{itemize}
\item To provide a secure and automated attendance process for classrooms.
\item To ensure that students are physically present near the teacher device while marking attendance.
\item To verify student identity using enrolled face data before attendance acceptance.
\item To generate teacher-side visibility through live and historical attendance views.
\item To prevent duplicate attendance entries for the same class slot.
\end{itemize}

\section{Problem Statement}
Current attendance systems in many institutions do not provide strong assurance that the student who marks attendance is both physically present and correctly identified. A single-factor approach can be bypassed. Therefore, the problem is to design and implement a practical multi-factor attendance workflow that is secure, scalable, and easy to operate in real classroom settings.

\section{Scope of the Project}
The current implementation supports:
\begin{itemize}
\item Teacher session initiation with class date and period.
\item Session-bound OTP generation and validation.
\item Bluetooth proximity verification for classroom presence.
\item Face verification flow integrated in student attendance marking.
\item Attendance storage in backend with timestamp and slot details.
\item Teacher dashboard modules for live list, history, and student reports.
\end{itemize}

The scope is focused on classroom-level operation in undergraduate institutions and can be extended further to department-wide analytics in future.

\section{Organization of the Report}
This report is organized into seven chapters. Chapter 1 introduces the project context and objectives. Chapter 2 presents the literature survey and comparative discussion. Chapter 3 explains system requirements and analysis. Chapter 4 describes design and architecture. Chapter 5 details implementation. Chapter 6 presents results and discussions. Chapter 7 concludes the work with future scope.

